% =============================================================================
% Seccion: Descomposicion matematica del problema de despacho.
% Fuente incluida desde Tesis/desarrollo.tex mediante Tesis/generados/.
% Mantener este archivo como fuente y sincronizar con `make sync` en Tesis/.
% =============================================================================
\section{Descomposición matemática del problema de despacho}
\label{sec:descomposicion_despacho}

Se reformula el despacho por agentes mediante un maestro de Dantzig--Wolfe y subproblemas locales, y se precisan las cotas de pricing necesarias para certificar su cierre. En el régimen lineal ensayado se reproduce el costo monolítico a la precisión reportada, con mayor tiempo de resolución. La línea permanece en desarrollo; la evaluación experimental y la extensión entera siguen abiertas.

La Sección~\ref{sec:marco_dw} presenta el criterio de elección por restricciones o variables complicantes~\cite{ConejoEtAl2006}; sus antecedentes hidrotérmicos se revisan en la Sección~\ref{sec:descomposicion-sthc} \cite{JimenezRedondoConejo1999,BorghettiEtAl2003,SifuentesVargas2007,ScuzziatoEtAl2018,ColonettiEtAl2022,SugishitaGrotheyMcKinnon2024}.

La estructura de la Sección~\ref{sec:analisis_estructural_modelo} separa Río Negro, Salto Grande y térmicas en bloques $x_i\in X_i$, con costo $C_i(x_i)$ y aporte horario $G_{i,t}(x_i)$. Este último incluye la participación nacional $r$ de Salto Grande. Los selectores $L_i$ agregan las generaciones locales; la falla $s_t=d_t-g_t\ge0$ equivale a limitar $g_t\le d_t$. Su costo $C_f(s)$ es la envolvente por tramos de~\eqref{eq:costo_falla_tramos}, representada por epígrafe. El problema acoplado es

\begin{align}
    \min_{x_1,\ldots,x_N,g,s}
    \quad
    &
    \sum_{i=1}^{N} C_i(x_i)
    +
    C_f(s)
    \label{eq:despacho_descomp_obj}
    \\
    \text{s.a.}
    \quad
    &
    x_i \in X_i,
    \qquad i=1,\ldots,N,
    \label{eq:despacho_descomp_local}
    \\
    &
    g_{i,t}=G_{i,t}(x_i),
    \qquad i=1,\ldots,N,\quad t=1,\ldots,T,
    \label{eq:despacho_descomp_generacion_local}
    \\
    &
    \sum_{i=1}^{N} g_{i,t} - g_t = 0,
    \qquad t=1,\ldots,T,
    \label{eq:despacho_descomp_acoplamiento}
    \\
    &
    g_t+s_t=d_t,
    \qquad t=1,\ldots,T,
    \label{eq:despacho_descomp_demanda}
    \\
    &
    s_t\geq 0,
    \qquad t=1,\ldots,T.
    \label{eq:despacho_descomp_falla}
\end{align}
La Figura~\ref{fig:dw_estructura_matriz} resume la estructura: bloques locales y filas globales que vinculan las trayectorias de generación.

\begin{figure}[htbp]
\centering
\begin{tikzpicture}[every node/.style={font=\footnotesize},
  blk/.style={ttproc, minimum size=11mm, inner sep=1pt},
  band/.style={minimum height=6.5mm, inner sep=2pt, align=center}]
  \node[blk] (a1) at (0,0)      {$A_1$};
  \node[blk] (a2) at (1.2,-1.2) {$A_2$};
  \node                at (2.1,-2.1) {$\ddots$};
  \node[blk] (aN) at (3.0,-3.0) {$A_N$};
  \draw[black!35] (-0.62,0.62) rectangle (3.62,-3.62);
  \node[ttres, band, minimum width=42.4mm, anchor=north] (cpl) at (1.5,-3.95)
     {acoplamiento\ \ $\textstyle\sum_{i} g_{i,t}-g_t=0$};
  \node[ttdata, band, minimum width=42.4mm, anchor=north] (dem) at (1.5,-4.9)
     {demanda/falla\ \ $g_t+s_t=d_t$};
  \draw[ttbrace] (3.82,0.55) -- (3.82,-3.55);
  \node[anchor=west, align=left, font=\scriptsize\itshape] at (4.0,-1.5)
     {bloques locales\\$X_i$ (separables)};
\end{tikzpicture}
\caption[Estructura por bloques de la matriz de restricciones]{Bloques locales separables $A_i$ y filas globales. Solo el balance $\sum_i g_{i,t}=g_t$ cruza todos los bloques; la demanda $g_t+s_t=d_t$ los coordina indirectamente mediante $g_t$. El ancho de las bandas es esquemático: se omiten las columnas de $g_t$ y $s_t$.}
\label{fig:dw_estructura_matriz}
\end{figure}

\subsection*{Criterio estructural y estrategias posibles}

El balance~\eqref{eq:despacho_descomp_acoplamiento} constituye la restricción complicante: al retirarlo, cada bloque conserva balances hidráulicos, rezagos y condiciones térmicas. Con multiplicadores $\lambda_t$, su relajación lagrangiana es

\begin{align}
    \mathcal{L}(x_1,\ldots,x_N,g,s,\lambda)
    =
    \sum_{i=1}^{N}
    \left[
        C_i(x_i)
        +
        \sum_{t=1}^{T}\lambda_t G_{i,t}(x_i)
    \right]
    -
    \sum_{t=1}^{T}\lambda_t g_t
    +
    C_f(s).
    \label{eq:lagrangiano_despacho}
\end{align}
Con multiplicadores fijos, la relajación se separa por agentes; sus soluciones pueden incumplir el balance y requieren recuperación primal.

Benders o Benders generalizado serían pertinentes si fijar decisiones globales dominantes, como trayectorias agregadas o decisiones estratégicas, simplificara los bloques. La descomposición por escenarios separaría realizaciones de incertidumbre y requeriría una formulación estocástica explícita; las instancias actuales son determinísticas o parametrizadas.


\subsection*{Estrategia adoptada}
\label{par:estrategia_descomposicion}

Se adopta Dantzig--Wolfe porque el acoplamiento por restricciones permite coordinar perfiles localmente factibles en un maestro. Es una elección estructural; la Tabla~\ref{tab:dw_lp_validacion} evalúa la reformulación lineal, sin comparar empíricamente las cuatro estrategias.

Para la trayectoria $G_i(x_i)=(G_{i,1}(x_i),\ldots,G_{i,T}(x_i))$, una columna $k$ es el par $(u_i^k,c_i^k)$ asociado a un perfil factible $x_i^k\in X_i$:
\begin{equation}
    u_i^k = G_i(x_i^k),
    \qquad
    c_i^k = C_i(x_i^k).
    \label{eq:columna_dw_definicion_general}
\end{equation}
La columna contiene generación $u_i^k$ [MW] y costo $c_i^k$ [USD]; conservar $x_i^k$ permite reconstruir la operación. Sea $\mathcal C_i$ el conjunto de columnas: $\mathcal P_i$ para el maestro completo o $\mathcal K_i^r\subseteq\mathcal P_i$ para el restringido en la iteración $r$. Ambos comparten la formulación

\begin{align}
    \min_{\lambda,s}
    \quad
    &
    \sum_{i=1}^{N}
    \sum_{k\in\mathcal C_i}
    c_i^k \lambda_{i,k}
    +
    C_f(s)
    \label{eq:maestro_dw_obj}
    \\
    \text{s.a.}
    \quad
    &
    \sum_{k\in\mathcal C_i}\lambda_{i,k}=1,
    \qquad i=1,\ldots,N,
    \label{eq:maestro_dw_convexidad}
    \\
    &
    \sum_{i=1}^{N}
    \sum_{k\in\mathcal C_i}
    u_{i,t}^{k}\lambda_{i,k}
    +
    s_t
    =
    d_t,
    \qquad t=1,\ldots,T,
    \label{eq:maestro_dw_demanda}
    \\
    &
    s_t\geq 0,
    \qquad t=1,\ldots,T,
    \label{eq:maestro_dw_falla}
    \\
    &
    \lambda_{i,k}\geq 0,
    \qquad i=1,\ldots,N,\quad k\in\mathcal C_i.
    \label{eq:maestro_dw_lambda}
\end{align}
El maestro combina perfiles para cubrir demanda y penalizar falla; $g_t=\sum_i\sum_{k\in\mathcal C_i}u_{i,t}^k\lambda_{i,k}$. Como $\mathcal P_i$ no se enumera, el algoritmo resuelve maestros restringidos y utiliza sus duales para generar columnas mediante~\eqref{eq:pricing_local_obj}--\eqref{eq:pricing_local_factibilidad}. Agrega perfiles de costo reducido negativo hasta cerrar el pricing o alcanzar otro criterio de parada.

La Figura~\ref{fig:dw_pipeline_despacho} muestra la inicialización, el intercambio de duales y columnas y las condiciones de parada.

\begin{figure}[htbp]
\centering
\begin{tikzpicture}[node distance=7mm and 8mm, every node/.style={font=\footnotesize}]
  \node[ttdata, text width=28mm] (seed)
    {\textbf{Inicialización} $\mathcal{K}_i^0$\\
     \scriptsize columnas a precios $\gamma\,\pi$\\
     \scriptsize \textbf{+ LP reparado}};
  \node[ttproc, text width=28mm, right=8mm of seed] (master)
    {\textbf{Maestro LP}\\
     \scriptsize acoplamiento y convexidad\\
     \scriptsize duales $(\pi_t,\mu_i)$};
  \node[ttproc, text width=28mm, right=8mm of master] (pricing)
    {\textbf{Pricing por agente}\\
     \scriptsize factibilidad local $X_i$\\
     \scriptsize columna $(u_i^k,c_i^k)$};
  \node[ttres, text width=28mm, below=10mm of master] (stop)
    {\textbf{Parada}\\
     \scriptsize certificada si\\
     \scriptsize $\ell_i-\mu_i\ge-\varepsilon_i\ \forall i$;\\
     \scriptsize o por estancamiento\\
     \scriptsize (heurístico)};
  \node[ttcore, text width=28mm, right=8mm of stop] (final)
    {\textbf{Maestro entero}\\
     \scriptsize opcional si hay\\
     \scriptsize variables enteras\\
     \scriptsize $\lambda_{i,k}\in\{0,1\}$};
  \draw[ttflow] (seed) -- (master);
  \draw[ttflow] ([yshift=2mm]master.east) -- ([yshift=2mm]pricing.west)
     node[midway, above, font=\scriptsize] {$\pi,\mu$};
  \draw[ttflow] ([yshift=-2mm]pricing.west) -- ([yshift=-2mm]master.east)
     node[midway, below=7mm, font=\scriptsize] {$\bar c_i^k<-\varepsilon$};
  \draw[ttflowdash] (master) -- (stop);
  \draw[ttflow] (stop) -- (final);
\end{tikzpicture}
\caption[Generación de columnas sobre el modelo de despacho]{Generación de columnas con precios iniciales escalados por un factor adimensional $\gamma$ y perfiles del LP reparado (Sección~\ref{subsec:dw_lp_disuelve}). Con duales óptimos del maestro original y cotas inferiores válidas $\ell_i$ del pricing, la parada certificada limita la brecha a $\sum_i\varepsilon_i$ [USD]; el estancamiento es heurístico. El maestro entero selecciona perfiles cuando los bloques conservan variables enteras; se omite bajo las hipótesis lineales de la Proposición~\ref{prop:dw_exacto_lineal}.}
\label{fig:dw_pipeline_despacho}
\end{figure}

Si los bloques conservan variables enteras, puede resolverse al final el maestro restringido con $\lambda_{i,k}\in\{0,1\}$, definido en~\eqref{eq:maestro_dw_binario}. Su alcance es el siguiente.

\begin{observacion}[Alcance del maestro entero restringido]
\label{obs:maestro_entero_restringido}
El cierre exacto del \emph{pricing} certifica el óptimo del maestro lineal completo.
Con bloques enteros, ese maestro optimiza sobre los conjuntos locales convexificados;
su valor es una cota inferior del MIP. El maestro entero restringido, si es factible,
selecciona una columna de cada conjunto generado y satisface el acoplamiento: entrega
una cota superior del MIP, óptima únicamente entre esas selecciones. No optimiza sobre
la envolvente convexa de la unión de columnas de bloques diferentes.
La igualdad de ambas cotas certifica optimalidad; para obtener garantías generales
puede recurrirse a una exploración completa de \emph{branch-and-price}, con reglas de
ramificación y \emph{pricing} compatibles \cite{DesrosiersLubbecke2005}.

La integralidad de cada poliedro local no garantiza la integralidad del maestro acoplado.
Si todos los bloques son continuos y convexos, en cambio, la reconstrucción mediante
combinaciones convexas es factible y el maestro entero deja de ser necesario
(Proposición~\ref{prop:dw_exacto_lineal}). Imponerlo sobre un conjunto finito de columnas
puede empeorar el objetivo incluso en ese caso lineal.
\end{observacion}

Se compara costo, falla, factibilidad, tiempo, columnas e iteraciones con la resolución directa de la formulación global, denominada \emph{monolítica}. El \emph{pricing} utiliza el solver; la asistencia mediante aprendizaje se estudia en la Sección~\ref{sec:hibridacion}.

\subsection{Construcci\'on de los subproblemas locales}
\label{subsec:construccion_subproblemas_locales}

Dados los precios del maestro $\pi=(\pi_1,\ldots,\pi_T)$, el agente $i$ genera un perfil factible resolviendo

\begin{align}
    \min_{x_i}
    \quad
    &
    C_i(x_i)
    -
    \sum_{t=1}^{T}\pi_t G_{i,t}(x_i)
    \label{eq:pricing_local_obj}
    \\
    \text{s.a.}
    \quad
    &
    x_i\in X_i.
    \label{eq:pricing_local_factibilidad}
\end{align}
El dual $\pi_t$ representa el valor marginal del aporte horario. La solución define $(u_i^k,c_i^k)$ mediante~\eqref{eq:columna_dw_definicion_general}; se conserva también $x_i^k$ para reconstruir el despacho. En los bloques hidráulicos, respetar máximos horarios no asegura balances, rezagos ni volúmenes. Por ello, las cotas

\begin{equation}
    \overline u_{i,t}
    =
    \max_{x_i\in X_i} G_{i,t}(x_i),
    \qquad t=1,\ldots,T,
    \label{eq:maximo_producible_periodo}
\end{equation}
son diagnósticas: los perfiles completos se generan sobre $X_i$. La Sección~\ref{sec:hibridacion} estudia propuestas aprendidas para este \emph{pricing}; la aproximación de la Sección~\ref{sec:aproximacion_redes_neuronales} propone directamente el despacho.

\subsection{Construcci\'on del problema maestro}
\label{subsec:construccion_problema_maestro}

En la iteración $r$, el maestro utiliza~\eqref{eq:maestro_dw_obj}--\eqref{eq:maestro_dw_lambda} con $\mathcal C_i=\mathcal K_i^r$. Sean $\mu_i$ [USD] y $\pi_t$ [USD/MW] los duales de convexidad y demanda. El costo reducido de una columna, en USD, es

\begin{equation}
    \bar c_i^k
    =
    c_i^k
    -
    \mu_i
    -
    \sum_{t=1}^{T}\pi_t u_{i,t}^{k}.
    \label{eq:costo_reducido_columna}
\end{equation}
El término $\mu_i$ es constante en el pricing. La columna se incorpora si $\bar c_i^k<-\varepsilon$, con tolerancia $\varepsilon>0$ [USD].

El maestro entero restringido de la Observación~\ref{obs:maestro_entero_restringido} añade

\begin{equation}
    \lambda_{i,k}\in\{0,1\},
    \qquad i=1,\ldots,N,\quad k\in\mathcal{K}_i^r.
    \label{eq:maestro_dw_binario}
\end{equation}
La convexidad~\eqref{eq:maestro_dw_convexidad} impone entonces una columna por agente. Esta selección conserva la factibilidad local, con el alcance de la Observación~\ref{obs:maestro_entero_restringido}.

\subsection{Límites de la descomposición del MIP}
\label{subsec:intento_binario_dw}

El primer ensayo conserva las binarias $\varphi_{i,t}^{(r)}$ de selección del tramo de altura $r$, en la hora $t$ de la central $i$.

La aplicación directa no alcanza una resolución satisfactoria en los grupos multiagente de mercado. Intervienen tres dificultades: reconstruir decisiones enteras, disponer de columnas suficientes y resolver el maestro dentro del presupuesto. La selección de una columna por agente recupera integralidad, pero puede omitir perfiles necesarios (Observación~\ref{obs:maestro_entero_restringido}). En \texttt{instance024}, los costos reportados del LP compacto y del MIP son $121{,}188$ y $232{,}973$ USD; su diferencia no determina la brecha DW, cuya cota completa puede ser más fuerte.

También se observan perfiles concentrados en horas de mayor precio y descenso lento del objetivo (\emph{tailing-off}). Incluso con bloques lineales pueden aparecer perfiles extremos y pivotes degenerados en el maestro: un costo reducido negativo no garantiza una mejora inmediata. Estos comportamientos no son exclusivos de las binarias, y un buen pricing local no asegura una buena selección entera en tiempo limitado.

\begin{figure}[H]
\centering
\includegraphics[width=0.78\textwidth]{figs/delta_obj_dw.pdf}
\caption[Diferencia de objetivo entre descomposición y MIP monolítico]{Diferencia relativa de objetivo de la descomposici\'on sobre el modelo \emph{binario}
frente al monolítico, por grupo: $\Delta z^{\mathrm{rel}}=(z_{\mathrm{DW}}-z_{\mathrm{MIP}})/|z_{\mathrm{MIP}}|$, cociente adimensional expresado en porcentaje en la figura ($100\Delta z^{\mathrm{rel}}$), con escala logarítmica simétrica. El Modelo~1 reproduce la referencia; los grupos multiagente de mercado concentran la brecha. Se distinguen integralidad, selección de columnas y convergencia.}
\label{fig:delta_obj_dw}
\end{figure}

Al inyectar las trayectorias del despacho de referencia, el maestro entero recupera su valor: esas columnas forman una selección factible y, si el despacho es óptimo, no pueden mejorarlo. El control identifica una limitación del conjunto generado, sin atribuir toda la convergencia a las binarias.

Se implementaron costo de falla con \emph{big-M} controlado, columnas iniciales obtenidas del LP, \emph{pricing} por etapas, filtrado mediante relajación, estabilización dual (Neame e \emph{in-out}) y separación entre cotas garantizadas y heurísticas \cite{DesrosiersLubbecke2005,PessoaSadykovUchoaVanderbeck2018,BulaichMehamdiLacroixMartin2024,MartinBauguionMagnoucheLeguay2024,RousseauGendreauFeillet2007}. También se ensayaron \emph{branch-and-price} externo y derivación al monolítico cuando el maestro entero selecciona falla.

Un bucle de cortes Dantzig--Wolfe en el espacio original reforzó el LP compacto en la raíz ($+24\%$ en \texttt{instance001}), pero su costo computacional degradó la cota final del Modelo~3.\footnote{Rama experimental no integrada en la línea principal.} La descomposición mejorada obtuvo una mediana de \emph{speedup} de $1{,}54$ en ese grupo, una mejora modesta que requiere ampliar la evaluación; en los restantes grupos el monolítico mantuvo la ventaja. La certificación entera requiere completar el \emph{branch-and-price}.

Un criterio preciso permite comparar las cotas. Si $X_i$ es el conjunto entero local y
$Q_i$ su relajación lineal, entonces $\operatorname{conv}(X_i)\subseteq Q_i$ y, para
minimización, $z_{\mathrm{LP}}\le z_{\mathrm{DW}}\le z_{\mathrm{IP}}$.
La igualdad $\operatorname{conv}(X_i)=Q_i$ en todos los bloques es suficiente para
$z_{\mathrm{DW}}=z_{\mathrm{LP}}$ \cite{DesrosiersLubbecke2005}; no es necesaria para
una instancia y un objetivo concretos, pues las dos formulaciones pueden compartir un
óptimo aunque sus conjuntos difieran. Si $z_{\mathrm{IP}}>z_{\mathrm{LP}}$, el cociente
$(z_{\mathrm{DW}}-z_{\mathrm{LP}})/(z_{\mathrm{IP}}-z_{\mathrm{LP}})$ mide la fracción
de brecha cerrada, siempre que se usen valores óptimos certificados de los tres problemas.
Una brecha pequeña para un objetivo tampoco demuestra integralidad del poliedro local.

Los regímenes ofrecen márgenes distintos. En el Modelo~1, la brecha del LP compacto es $\sim0{,}015\%$ (\texttt{caso01a}: $\sim5\times10^4$ USD sobre $\sim3{,}4\times10^8$ USD), por lo que queda escaso margen para fortalecer la cota. En el Modelo~3, \texttt{instance001} registra $z_{\rm LP}=98{,}410$ USD frente a un incumbente de $230{,}516$ USD ($57\%$ por debajo); el monolítico conserva una brecha de $43{,}5\%$ tras $240\,\mathrm{s}$. El \emph{branch-and-price} en la raíz consumió $571\,\mathrm{s}$ y $21$ iteraciones, con $99{,}8\%$ del tiempo en \emph{pricing} entero de la central no lineal, sin cerrar la cota DW. Esta corrida muestra que el costo de obtener una cota más fuerte puede absorber su ventaja potencial.

El costo del \emph{pricing} entero y la selección de columnas motivan la evaluación sobre cabeza congelada. La comparación cambia convexidad, coeficientes, conjuntos factibles de producción e inicialización: diagnostica un régimen distinto, sin aislar el efecto de retirar las binarias.

\subsection{Exactitud de la reformulación sobre la aproximación lineal}
\label{subsec:dw_lp_disuelve}

Esta campaña obtiene la cabeza congelada mediante iteración de Picard, sin cargar una solución MIP; el posprocesamiento anclado emplea los promedios del despacho MIP (Sección~\ref{sec:exp-cabeza-congelada}). Cambiar promedios o tangentes puede producir LP distintos. La proposición siguiente demuestra la exactitud de Dantzig--Wolfe para cada LP fijado y permite reconstruir sus soluciones sin seleccionar perfiles binarios; no establece equivalencia con la producción física no lineal. Las columnas disponibles y el cierre del pricing siguen determinando la calidad del maestro restringido.

\begin{proposicion}[Exactitud de Dantzig--Wolfe sobre bloques lineales]
\label{prop:dw_exacto_lineal}
Fíjense la cabeza y todos los coeficientes del problema acoplado
\eqref{eq:despacho_descomp_obj}--\eqref{eq:despacho_descomp_falla}.
Supóngase que cada $X_i$ es un poliedro no vacío y acotado, que $C_i$ y $G_i$ son afines,
y que el costo de falla $C_f$ tiene una representación lineal mediante epígrafe.
Supóngase además que el problema acoplado es factible y tiene óptimo finito. Entonces:
\begin{enumerate}
 \item[\textnormal{(a)}] toda combinación convexa $\hat x_i=\sum_k\lambda_{i,k}x_i^k$
 de columnas factibles pertenece a $X_i$;
 \item[\textnormal{(b)}] el maestro lineal completo, con todos los vértices de cada $X_i$,
 tiene el mismo valor óptimo que el problema acoplado;
 \item[\textnormal{(c)}] cualquier solución factible de un maestro restringido se reconstruye
 mediante esos $\hat x_i$ en una solución factible del problema acoplado, al mismo costo.
 Esa solución es óptima para el problema acoplado si también es óptima para el maestro
 completo; el cierre exacto del \emph{pricing} proporciona ese certificado.
\end{enumerate}
\end{proposicion}
\begin{proof}
La convexidad de $X_i$ prueba (a). Por Minkowski--Weyl, cada politopo es la envolvente
convexa de sus vértices. Por tanto, todo punto factible del problema acoplado admite la
representación $x_i=\sum_k\lambda_{i,k}x_i^k$ y genera una solución factible del maestro
completo. Recíprocamente, (a) reconstruye los bloques de cualquier solución del maestro.
La afinidad y $\sum_k\lambda_{i,k}=1$ dan
\[
 C_i(\hat x_i)=\sum_k\lambda_{i,k}C_i(x_i^k),\qquad
 G_i(\hat x_i)=\sum_k\lambda_{i,k}G_i(x_i^k).
\]
Conservar la misma falla $s$ y su epígrafe mantiene el acoplamiento y el objetivo en
ambas direcciones. Esto prueba (b) y la primera afirmación de (c), que también vale
para un subconjunto de columnas. La afirmación de optimalidad de (c) sigue de (b) y de
la Proposición~\ref{prop:cierre_pricing_dw}; las variables explícitas de falla conservan
sus restricciones duales, mientras el \emph{pricing} verifica las columnas ausentes.
\end{proof}

La acotación debe comprobarse para todas las variables locales, incluidas las auxiliares;
si hay direcciones de recesión relevantes, deben incorporarse columnas de rayos o utilizar
una formulación equivalente acotada. La reconstrucción es factible para el LP fijado:
la comprobación de producción física corresponde a la auditoría de cabeza congelada.

El resultado no dice que imponer $\lambda\in\{0,1\}$ conserve el óptimo. Por ejemplo,
con un bloque $X=[0,1]$, generación $G(x)=x$, costo local cero, demanda $d=1/2$ y costo
de falla $C_f(s)=s$, las columnas $0$ y $1$ bastan para el maestro lineal completo.
La mezcla $(1/2,1/2)$ reconstruye $x=1/2$ con costo cero. El maestro binario sobre esas
mismas columnas debe elegir $x=0$, pues $x=1$ excede la demanda; su costo es $1/2$.
Incluir la columna agregada $x=1/2$ elimina esa brecha. Así, una brecha de integerización
nula en una campaña describe el conjunto de columnas generado, no una consecuencia
general de que los bloques sean lineales.

\paragraph{Cierre, tolerancia e inicialización.}
El certificado de~\eqref{eq:dw-cierre-tolerancia} usa mínimos globales o cotas inferiores válidas de todos los pricing, con los duales originales del maestro. Las tolerancias $\varepsilon_i$ [USD] se acumulan; el estancamiento no certifica optimalidad.

En las corridas, una columna inicial obtenida de la aproximación lineal y reparada como perfil
balanceado mejora el comportamiento dentro del presupuesto disponible. No constituye
por sí sola un teorema de convergencia. Si se inyectan las trayectorias de un óptimo del
mismo LP, el maestro restringido contiene ese óptimo desde el inicio; si sólo se inyecta
una solución factible reparada, esa conclusión requiere una comprobación adicional.
El \emph{tailing-off} puede persistir sobre bloques lineales: en \texttt{instance024}
se reporta un objetivo aproximadamente $240\%$ superior al óptimo tras $200$
iteraciones sin la columna inicial informada.

\begin{figure}[htbp]
\centering
\includegraphics[width=0.96\textwidth]{figs/convergencia_dw.pdf}
\caption[Convergencia de generación de columnas sobre el LP]{Evolución de la generación de columnas sobre cabeza congelada con columna inicial
lineal reparada. Se muestran las medianas del objetivo del maestro restringido y del
menor costo reducido por iteración. El objetivo no aumenta cuando se agregan columnas
sin retirar las existentes y se resuelve el mismo maestro. Una mediana de costos
reducidos cercana a cero no certifica el cierre en todas las instancias: el control
requiere cotas de \emph{pricing} por bloque e instancia y la tolerancia acumulada.}
\label{fig:dw_convergencia_lp}
\end{figure}

\subsection{Comparación numérica sobre las 306 instancias}
\label{subsec:dw_lp_resultados}

Con columna inicial reparada y parada por estancamiento, DW reproduce el objetivo del LP monolítico dentro del umbral de comparación en las $306$ instancias (Tabla~\ref{tab:dw_lp_validacion}). La brecha de selección entera también es numéricamente nula, pero depende de las columnas generadas. La coincidencia observada concuerda con la reformulación exacta; no sustituye los certificados primal y dual por instancia.

\begin{table}[H]
\centering
\caption[Validación de Dantzig--Wolfe sobre las 306 instancias lineales]{DW frente al monolítico de cabeza congelada, en nodo dedicado del clúster. «DW $\approx$ Mono» cuenta los casos con $|\Delta z^{\rm rel}|<10^{-6}$; la brecha entera es $(z_{\rm ent}-z_{\rm LP})/|z_{\rm LP}|$. Ambos cocientes son adimensionales. El tiempo se expresa como mediana (IQR), en segundos.}
\label{tab:dw_lp_validacion}
\small
\renewcommand{\arraystretch}{1.2}
\begin{tabularx}{\textwidth}{@{}>{\raggedright\arraybackslash}Xccccc@{}}
\toprule
\textbf{Modelo} & $n$ & DW $\approx$ Mono & $\mathrm{med}\,|\Delta z^{\mathrm{rel}}|$ & \shortstack{Brecha entera\\(máximo)} & $T_{\mathrm{DW}}$ (s) \\
\midrule
Modelo~1   & $204$ & $204/204$ & $9\cdot10^{-16}$ & $0$              & $1{,}27$ ($0{,}28$) \\
Modelo~2      & $68$  & $68/68$   & $1\cdot10^{-13}$ & $9\cdot10^{-14}$ & $0{,}91$ ($0{,}30$) \\
Modelo~3 & $34$  & $34/34$   & $2\cdot10^{-13}$ & $0$              & $2{,}34$ ($0{,}36$) \\
\bottomrule
\end{tabularx}
\end{table}

Las instancias \texttt{instance024} (Modelo~3) e \texttt{instance071} (Modelo~2), antes dominadas por falla, coinciden con la referencia lineal. La comparación cambia tanto la aproximación de producción como la inicialización; no aísla el efecto de la convexidad.

$T_{\rm DW}$ incluye maestro y pricing. El monolítico resuelve en fracciones de segundo a pocos segundos y resulta comparable o más rápido; no se registra aquí su tiempo por grupo. El aporte es reconstruir los despachos de los LP aproximados y disponer de una base para generar columnas aprendidas, con costo adicional de coordinación.


\paragraph{Certificación por pasada de pricing.}
La implementación revisada calcula
$z_{\rm MR}+\sum_i\min\{0,\ell_i-\mu_i\}$ con cotas inferiores del solver
para todos los bloques de una pasada pura. Una solución incumbente sólo se usa
para proponer columnas. Los duales originales del maestro se conservan intactos;
su recorte se aplica únicamente a los precios heurísticos. Las pasadas con
maestro perturbado no certifican el problema original. El registro distingue
la ausencia de columnas nuevas de un cierre acreditado y guarda las cotas de
cada bloque con precisión suficiente para reproducir la suma. Una cota válida
anterior permanece válida aunque un pricing posterior se interrumpa.
La comparación pequeña con el LP monolítico del mismo caso y las pruebas de
pricing truncado verifican estas distinciones. Los certificados siguen sujetos
a las tolerancias numéricas del solver y no certifican la producción no convexa.


\begin{figure}[htbp]
\centering
\includegraphics[width=\textwidth]{figs/rev_certificados_computables.pdf}
\caption[Certificados computables de descomposición y reparación]{Certificados reproducidos en la revisión. Arriba: caso sintético SG--Ter--Mer
 de ocho horas; el objetivo del maestro coincide pronto con el LP monolítico,
 mientras la cota inferior mejora hasta cerrar el pricing. Abajo: tres escenarios
 de Río Negro de 360 horas, con la reparación cuadrática corregida. La referencia
 $w$ proviene del LP económico; $D(w)$ no se identifica con $D^\star$. El
 certificado numérico agrega el residuo de primer orden y el error de referencia
 de la Proposición~\ref{prop:red-cota-computable}; la línea muestra el factor del
 caso ideal. Estas corridas no representan las campañas históricas completas.}
\label{fig:certificados-computables}
\end{figure}

\subsection{Conclusi\'on: el rol de la descomposici\'on}
\label{subsec:conclusion_descomposicion}

La reformulación reproduce los costos de los LP ensayados sin acelerar su resolución. La extensión entera debe abordar el costo del pricing, la degeneración y la selección de perfiles; las corridas no aíslan a las binarias como causa exclusiva. Las ventajas en escalas mayores, escenarios o nuevos componentes discretos requieren evaluación. La Sección~\ref{sec:hibridacion} estudia propuestas de columnas aprendidas, sujetas a validación y al cierre del pricing.

Los perfiles pueden reutilizarse si permanecen iguales todos los parámetros locales: aportes, volumen inicial, valor del agua, cabeza, tangentes y restantes restricciones. Cambiar sólo demanda y costo de falla modifica el maestro, pero exige verificar de nuevo el \emph{pricing} para certificar optimalidad. Los duales describen valores marginales del LP con su orientación de filas; pueden no ser únicos ante degeneración. Estas propiedades sustentan resolución repetida y coordinación modular, sin establecer una ventaja temporal sobre el solver directo.
